\documentclass[10pt,a4paper]{amsart}
\usepackage[margin=19mm]{geometry}
\usepackage{amsmath,amssymb,mathtools}
\usepackage[scr=rsfs,cal=euler]{mathalfa}
\usepackage{xfrac}
\usepackage[unicode,hidelinks,bookmarks=false]{hyperref}
\newcommand{\ie}{i.e.\ }
\newcommand{\cf}{cf.\ }
\newcommand{\resp}{respectively\ }
\newcommand{\Subsection}{\subsection}
\providecommand{\nobreakdash}{\nobreak\hbox{-}}
\setlength{\emergencystretch}{2em}
\multlinegap0pt
\usepackage{amssymb}
\newtheorem{lemma}{Lemma}
\newtheorem{proposition}{Proposition}
\title{\textsf{Monotone convergence theorems equivalent to Markov's principle}}
\author{Douglas S. Bridges}
\date{Source: arXiv:2606.08863v1 (CC BY 4.0)}
\begin{document}
\maketitle

\noindent\emph{Source and licence.} \textsf{Monotone convergence theorems equivalent to Markov's principle}. Version arXiv:2606.08863v1 (CC BY 4.0), distributed by arXiv under the Creative Commons Attribution 4.0 International licence, http://creativecommons.org/licenses/by/4.0/, as recorded in the arXiv licence metadata for this exact version and shown on the abstract page https://arxiv.org/abs/2606.08863v1. Full licence notice: \texttt{licenses/arxiv-2606.08863-CC-BY-4.0.txt}.\par\smallskip

\noindent\textit{Editorial scope.} Counted unit: the article's proposition 2 (source0.tex lines 207--216). The article's complete original proof of it is delivered with it (source0.tex lines 218--267, one \texttt{proof} environment of 50 lines). No mathematics is written, repaired or re-derived: the statement and the proof are the article's own lines, in the article's own notation, and no retained line is substituted or re-flowed.  Retained uncounted as the source-local context the proof needs: lines 191--195.  Nothing was omitted from any retained span, so the package carries no omission record.  No retained line contains a citation, so the wrapper carries no bibliography and no bibliography entry.  No retained line contains a figure environment, an image-inclusion command or drawing code, so no image is delivered and the package declares no asset dependency.  Editorial formatting only, in addition: an AMS \texttt{amsart} preamble that restates the article's own declarations and macros used by the retained lines, namely the environments \texttt{lemma}, \texttt{proposition} on separate counters, together with the standard packages those lines need.  The retained environments print in this document as Lemma 1, Proposition 1 in order of appearance, whereas the article prints the same items with the numbering of its own full document.  The complete native source of the article as distributed in the arXiv e-print for this exact version is bundled unchanged as \texttt{sources/202364/source0.tex}.
\begin{lemma}
\label{130511a3}Let $x\in\left[  0,1\right]  $ satisfy $\lnot\left(
x=1\right)  $. Then the sequence $\left(  x^{n}\right)  _{n\geqslant1}$ is
provisionally convergent to $0$. If it converges to $0$, then $x<1$.
\end{lemma}

\begin{proposition}
\label{2}The following are equivalent over\emph{\ \textbf{BISH}}:

\begin{itemize}
\item[\emph{(i)}] \emph{\textbf{MP}}.

\item[\emph{(ii)}] Every decreasing sequence of nonnegative real numbers that
is provisionally convergent to $0$ actually converges to $0$.
\end{itemize}
\end{proposition}

\begin{proof}
Assume \textbf{MP}. Let $\mathbf{x\equiv}\left(  x_{n}\right)  _{n\geqslant1}
$ be a decreasing sequence of nonnegative numbers converging provisionally to
$0$. Given $\varepsilon>0$, assume that $x_{n}>\varepsilon$ for all $n$.
Setting $n_{0}=1$, construct an increasing binary sequence $\left(
\lambda_{n}\right)  _{n\geqslant1}$ such that%
\begin{align*}
\lambda_{n}=0  & \Rightarrow\forall_{k\leqslant n}\left(  \left\vert
x_{k}-x_{1}\right\vert <\varepsilon\right)  ,\\
\lambda_{n}=1-\lambda_{n-1}  & \Rightarrow\left\vert x_{n}-x_{1}\right\vert
>\varepsilon/2.
\end{align*}
Since $\mathbf{x}$ is provisionally convergent to $0$, if $\lambda_{n}=0$ for
all $n$, then $\left\vert x_{1}\right\vert <\varepsilon$, which is absurd. It
follows from \textbf{MP} that there exists $n_{1}$ such that $\left\vert
x_{n_{1}}-x_{1}\right\vert >\varepsilon/2$. Repeating this with $x_{n_{1}}$
replacing $x_{1}$, and with the sequence $\left(  x_{n_{1}+1},x_{n_{1}%
+2},\ldots\right)  $ replacing $\mathbf{x}$, we obtain $n_{2}>n_{1}$ such that
$\left\vert x_{n_{2}}-x_{n_{1}}\right\vert >\varepsilon/2$. Carrying on in
this way, we construct a strictly increasing sequence $\left(  n_{k}\right)
_{k\geqslant1}$ such that $\left\vert x_{n_{k}}-x_{n_{k-1}}\right\vert
>\varepsilon/2$ for each $k\geqslant1$. Choosing a positive integer%
\[
K\geqslant2\left(  1+\frac{x_{1}}{\varepsilon}\right)  ,
\]
since the sequence $\mathbf{x}$ is decreasing we now see that%
\begin{align*}
x_{n_{K}}  & =x_{n_{1}}+\sum_{k=2}^{K}\left(  x_{n_{k}}-x_{n_{k-1}}\right) \\
& \leqslant x_{1}-\sum_{k=2}^{K}\left\vert x_{n_{k}}-x_{n_{k-1}}\right\vert
<x_{1}-\left(  K-2\right)  \frac{\varepsilon}{2}\leqslant0,
\end{align*}
a contradiction from which we conclude that%
\[
\lnot\forall_{n}\left(  x_{n}>\varepsilon\right)  .
\]
Applying \textbf{MP} once more, we deduce that there exists $\nu$ such that
$x_{\nu}<2\varepsilon$. Hence $0\leqslant x_{n}\leqslant x_{\nu}<2\varepsilon$
for all $n\geqslant\nu$. Since $\varepsilon>0$ is arbitrary, the proof that
(i) implies (ii) is complete.

Now assume (ii), and consider $a\in\left[  0,1\right]  $ with $\lnot\left(
a=0\right)  .$ By Lemma \ref{130511a3}, the sequence $\left(  \left(
1-a\right)  ^{n}\right)  _{n\geqslant1}$ is provisionally convergent to $0$,
and if it converges to $0$, then $1-a<1$ and therefore $a\neq0$. Thus the
statement in question implies that%
\[
\forall_{x\in\mathbb{R}}\left(  \lnot(x=0)\Rightarrow x\neq0\right)  ,
\]
which is equivalent to \textbf{MP}. Hence (ii) implies (i).
\end{proof}%

\end{document}
